% Options for packages loaded elsewhere
% Options for packages loaded elsewhere
\PassOptionsToPackage{unicode}{hyperref}
\PassOptionsToPackage{hyphens}{url}
\PassOptionsToPackage{dvipsnames,svgnames,x11names}{xcolor}
%
\documentclass[
  11pt,
  letterpaper,
]{article}
\usepackage{xcolor}
\usepackage[top=1in,bottom=1in,left=1in,right=1in]{geometry}
\usepackage{amsmath,amssymb}
\setcounter{secnumdepth}{5}
\usepackage{iftex}
\ifPDFTeX
  \usepackage[T1]{fontenc}
  \usepackage[utf8]{inputenc}
  \usepackage{textcomp} % provide euro and other symbols
\else % if luatex or xetex
  \usepackage{unicode-math} % this also loads fontspec
  \defaultfontfeatures{Scale=MatchLowercase}
  \defaultfontfeatures[\rmfamily]{Ligatures=TeX,Scale=1}
\fi
\usepackage{lmodern}
\ifPDFTeX\else
  % xetex/luatex font selection
\fi
% Use upquote if available, for straight quotes in verbatim environments
\IfFileExists{upquote.sty}{\usepackage{upquote}}{}
\IfFileExists{microtype.sty}{% use microtype if available
  \usepackage[]{microtype}
  \UseMicrotypeSet[protrusion]{basicmath} % disable protrusion for tt fonts
}{}
\makeatletter
\@ifundefined{KOMAClassName}{% if non-KOMA class
  \IfFileExists{parskip.sty}{%
    \usepackage{parskip}
  }{% else
    \setlength{\parindent}{0pt}
    \setlength{\parskip}{6pt plus 2pt minus 1pt}}
}{% if KOMA class
  \KOMAoptions{parskip=half}}
\makeatother
% Make \paragraph and \subparagraph free-standing
\makeatletter
\ifx\paragraph\undefined\else
  \let\oldparagraph\paragraph
  \renewcommand{\paragraph}{
    \@ifstar
      \xxxParagraphStar
      \xxxParagraphNoStar
  }
  \newcommand{\xxxParagraphStar}[1]{\oldparagraph*{#1}\mbox{}}
  \newcommand{\xxxParagraphNoStar}[1]{\oldparagraph{#1}\mbox{}}
\fi
\ifx\subparagraph\undefined\else
  \let\oldsubparagraph\subparagraph
  \renewcommand{\subparagraph}{
    \@ifstar
      \xxxSubParagraphStar
      \xxxSubParagraphNoStar
  }
  \newcommand{\xxxSubParagraphStar}[1]{\oldsubparagraph*{#1}\mbox{}}
  \newcommand{\xxxSubParagraphNoStar}[1]{\oldsubparagraph{#1}\mbox{}}
\fi
\makeatother


\usepackage{longtable,booktabs,array}
\usepackage{calc} % for calculating minipage widths
% Correct order of tables after \paragraph or \subparagraph
\usepackage{etoolbox}
\makeatletter
\patchcmd\longtable{\par}{\if@noskipsec\mbox{}\fi\par}{}{}
\makeatother
% Allow footnotes in longtable head/foot
\IfFileExists{footnotehyper.sty}{\usepackage{footnotehyper}}{\usepackage{footnote}}
\makesavenoteenv{longtable}
\usepackage{graphicx}
\makeatletter
\newsavebox\pandoc@box
\newcommand*\pandocbounded[1]{% scales image to fit in text height/width
  \sbox\pandoc@box{#1}%
  \Gscale@div\@tempa{\textheight}{\dimexpr\ht\pandoc@box+\dp\pandoc@box\relax}%
  \Gscale@div\@tempb{\linewidth}{\wd\pandoc@box}%
  \ifdim\@tempb\p@<\@tempa\p@\let\@tempa\@tempb\fi% select the smaller of both
  \ifdim\@tempa\p@<\p@\scalebox{\@tempa}{\usebox\pandoc@box}%
  \else\usebox{\pandoc@box}%
  \fi%
}
% Set default figure placement to htbp
\def\fps@figure{htbp}
\makeatother





\setlength{\emergencystretch}{3em} % prevent overfull lines

\providecommand{\tightlist}{%
  \setlength{\itemsep}{0pt}\setlength{\parskip}{0pt}}



 


\makeatletter
\@ifpackageloaded{tcolorbox}{}{\usepackage[skins,breakable]{tcolorbox}}
\@ifpackageloaded{fontawesome5}{}{\usepackage{fontawesome5}}
\definecolor{quarto-callout-color}{HTML}{909090}
\definecolor{quarto-callout-note-color}{HTML}{0758E5}
\definecolor{quarto-callout-important-color}{HTML}{CC1914}
\definecolor{quarto-callout-warning-color}{HTML}{EB9113}
\definecolor{quarto-callout-tip-color}{HTML}{00A047}
\definecolor{quarto-callout-caution-color}{HTML}{FC5300}
\definecolor{quarto-callout-color-frame}{HTML}{acacac}
\definecolor{quarto-callout-note-color-frame}{HTML}{4582ec}
\definecolor{quarto-callout-important-color-frame}{HTML}{d9534f}
\definecolor{quarto-callout-warning-color-frame}{HTML}{f0ad4e}
\definecolor{quarto-callout-tip-color-frame}{HTML}{02b875}
\definecolor{quarto-callout-caution-color-frame}{HTML}{fd7e14}
\makeatother
\makeatletter
\@ifpackageloaded{caption}{}{\usepackage{caption}}
\AtBeginDocument{%
\ifdefined\contentsname
  \renewcommand*\contentsname{Table of contents}
\else
  \newcommand\contentsname{Table of contents}
\fi
\ifdefined\listfigurename
  \renewcommand*\listfigurename{List of Figures}
\else
  \newcommand\listfigurename{List of Figures}
\fi
\ifdefined\listtablename
  \renewcommand*\listtablename{List of Tables}
\else
  \newcommand\listtablename{List of Tables}
\fi
\ifdefined\figurename
  \renewcommand*\figurename{Figure}
\else
  \newcommand\figurename{Figure}
\fi
\ifdefined\tablename
  \renewcommand*\tablename{Table}
\else
  \newcommand\tablename{Table}
\fi
}
\@ifpackageloaded{float}{}{\usepackage{float}}
\floatstyle{ruled}
\@ifundefined{c@chapter}{\newfloat{codelisting}{h}{lop}}{\newfloat{codelisting}{h}{lop}[chapter]}
\floatname{codelisting}{Listing}
\newcommand*\listoflistings{\listof{codelisting}{List of Listings}}
\makeatother
\makeatletter
\makeatother
\makeatletter
\@ifpackageloaded{caption}{}{\usepackage{caption}}
\@ifpackageloaded{subcaption}{}{\usepackage{subcaption}}
\makeatother
\usepackage{bookmark}
\IfFileExists{xurl.sty}{\usepackage{xurl}}{} % add URL line breaks if available
\urlstyle{same}
\hypersetup{
  pdftitle={Field Safety},
  pdfauthor={Marc Los Huertos},
  colorlinks=true,
  linkcolor={blue},
  filecolor={Maroon},
  citecolor={Blue},
  urlcolor={Blue},
  pdfcreator={LaTeX via pandoc}}


\title{Field Safety}
\author{Marc Los Huertos}
\date{2026-05-13}
\begin{document}
\maketitle

\renewcommand*\contentsname{Table of contents}
{
\hypersetup{linkcolor=}
\setcounter{tocdepth}{3}
\tableofcontents
}

\section{1. Purpose}\label{purpose}

This SOP covers field-based research activities and class field trips,
addressing transportation, weather-based, and biotic hazards. Although
we cannot predict all risks and hazards, it is everyone's responsibility
to identify and mitigate potential risks associated with field work.

\section{2. Scope}\label{scope}

\begin{itemize}
\tightlist
\item
  \textbf{Applies to:} All field-based research activities, class field
  trips, and off-campus data collection
\item
  \textbf{Does not apply to:} Laboratory-only procedures (see relevant
  Lab SOPs)
\item
  \textbf{Location:} All field sites used by the EA Program
\end{itemize}

\section{3. Responsibilities}\label{responsibilities}

\begin{longtable}[]{@{}
  >{\raggedright\arraybackslash}p{(\linewidth - 2\tabcolsep) * \real{0.2500}}
  >{\raggedright\arraybackslash}p{(\linewidth - 2\tabcolsep) * \real{0.7500}}@{}}
\toprule\noalign{}
\begin{minipage}[b]{\linewidth}\raggedright
Role
\end{minipage} & \begin{minipage}[b]{\linewidth}\raggedright
Responsibility
\end{minipage} \\
\midrule\noalign{}
\endhead
\bottomrule\noalign{}
\endlastfoot
Instructor / PI & Approve field plans, ensure training compliance \\
Trained Personnel & Follow this SOP, report deviations and hazards \\
New Students & Complete training and document review before field
work \\
\end{longtable}

Anyone conducting research in the field is required to have documented
training before commencing field work. Learning about risks and taking
the precautions described here is the first step to reducing the risks
associated with field work. Anyone working in the field should identify
potential risks and attempt to mitigate them.

\section{4. Safety and Hazards}\label{safety-and-hazards}

\begin{tcolorbox}[enhanced jigsaw, colback=white, titlerule=0mm, leftrule=.75mm, left=2mm, toptitle=1mm, rightrule=.15mm, title=\textcolor{quarto-callout-warning-color}{\faExclamationTriangle}\hspace{0.5em}{General Field Hazards}, opacitybacktitle=0.6, bottomrule=.15mm, breakable, opacityback=0, colframe=quarto-callout-warning-color-frame, coltitle=black, arc=.35mm, bottomtitle=1mm, colbacktitle=quarto-callout-warning-color!10!white, toprule=.15mm]

Reducing risks is one of our main priorities. Like many endeavors in
life, field work has certain risks. However, we can reduce these risks
dramatically by understanding our physical limits and how these can be
tested by uncontrolled environmental factors.

We rely on each individual to be mindful of their own limitations and to
advocate for themselves and others to reduce risk exposure.

\end{tcolorbox}

\subsection{Required Materials}\label{required-materials}

\begin{itemize}
\tightlist
\item
  First Aid Kit (kits are located in EA vehicles; see EA vehicle
  checklist to document their presence)
\item
  Cell phone(s)
\item
  Walkie-talkies (optional)
\item
  Personal Protective Equipment (PPE)
\end{itemize}

\subsection{Estimated Time}\label{estimated-time}

Preparing to go in the field can take a few minutes to several hours of
planning and gathering resources. A good rule of thumb is to double the
time you think it will take to get ready.

\section{5. Transportation Safety}\label{transportation-safety}

\subsection{5.1 Vehicle Requirements}\label{vehicle-requirements}

\begin{itemize}
\tightlist
\item
  Drivers must have valid driver's licenses and insurance
\item
  Drivers using EA vehicles must be approved by the college; forms can
  be found in the EA Department
\end{itemize}

\subsection{5.2 Planning}\label{planning}

\begin{itemize}
\tightlist
\item
  Identify meet-up and departure times and locations
\item
  Designate a driver
\item
  All state and local laws must be obeyed
\end{itemize}

\subsection{5.3 Responding to Vehicle
Emergencies}\label{responding-to-vehicle-emergencies}

Refer to the college's vehicle emergency procedures. Ensure all
occupants are safe, call emergency services if needed, and report the
incident to the EA Department.

\section{6. Weather-Related Hazards}\label{weather-related-hazards}

\subsection{6.1 Heat Exhaustion and
Stroke}\label{heat-exhaustion-and-stroke}

\begin{tcolorbox}[enhanced jigsaw, colback=white, titlerule=0mm, leftrule=.75mm, left=2mm, toptitle=1mm, rightrule=.15mm, title=\textcolor{quarto-callout-warning-color}{\faExclamationTriangle}\hspace{0.5em}{Risk Factors}, opacitybacktitle=0.6, bottomrule=.15mm, breakable, opacityback=0, colframe=quarto-callout-warning-color-frame, coltitle=black, arc=.35mm, bottomtitle=1mm, colbacktitle=quarto-callout-warning-color!10!white, toprule=.15mm]

\begin{itemize}
\tightlist
\item
  Working outside of shade in desert or arid climates, especially during
  hot times of day
\item
  Lack of shade coverage or proper sun protection
\item
  Dehydration
\end{itemize}

\end{tcolorbox}

\textbf{Prevention:}

\begin{itemize}
\tightlist
\item
  Use and reapply sunscreen
\item
  Drink plenty of water
\item
  Wear hats and other protective clothing
\end{itemize}

\textbf{Response:}

\begin{itemize}
\tightlist
\item
  Hydrate the affected person
\item
  Place ice packs on neck, armpits, etc.
\item
  If conditions worsen, leave the field and find air-conditioned
  conditions
\end{itemize}

\subsection{6.2 Cold Conditions}\label{cold-conditions}

\textbf{Risk factors:} Lack of sun, time of year, wind chill

\textbf{Prevention:} Proper layering and clothing; stay hydrated

\textbf{Response:} Eat, evacuate to warmth, get a professional
evaluation

\subsection{6.3 Lightning}\label{lightning}

\textbf{Prevention:} Check weather forecast; learn to identify
cumulonimbus clouds

\textbf{Response:} Assume lightning position and seek shelter
immediately; avoid open areas, tall isolated trees, and bodies of water

\subsection{6.4 Air Quality and
Pollution}\label{air-quality-and-pollution}

Working outside when ozone levels or other pollutants are elevated has a
detrimental effect on health, particularly for individuals with asthma
or other respiratory conditions. Hot, dry days worsen pollution, and
nearby forest fires degrade air quality.

Students with any respiratory disease should be cautious and report the
condition to their instructor or laboratory supervisor.

\textbf{Prevention:} Check air quality index (AQI) before field work;
postpone if AQI is unhealthy

\textbf{Response:} Leave the field if symptoms develop; seek medical
attention if needed

\section{7. Geological and Environmental
Hazards}\label{geological-and-environmental-hazards}

\subsection{7.1 Earthquake}\label{earthquake}

Tectonic movement may occur anywhere in California.

\textbf{Prevention:} Be aware of surroundings at all times

\textbf{Response:} Remain low to the ground; stay away from external
building frames; avoid sinkholes, gas lines, and utility wires

\subsection{7.2 Flooding}\label{flooding}

\textbf{Risk factors:} Proximity to bodies of water, excessive rainfall

\textbf{Prevention:} Check the forecast for rainfall; be aware of
surroundings

\textbf{Response:} Find the highest point in the area; do not attempt to
cross rushing water

\subsection{7.3 Altitude Illness
(HAPE/HACE)}\label{altitude-illness-hapehace}

\textbf{Risk factors:} High altitudes to which students have not had
time to acclimate; reduced oxygen

\textbf{Symptoms:} Headache, nausea, irritability

\textbf{Response:} If mild symptoms, acclimate and take it easy;
hydrate. If persistent vomiting and severe headache, evacuate to lower
altitude immediately.

\subsection{7.4 Human Debris and Waste}\label{human-debris-and-waste}

Illegal dumping of waste and leftover construction materials (sharp
metallic objects, etc.) bring risk of injury and exposure to long-term
illnesses like tetanus.

\textbf{Prevention:} Avoid and report; be aware of surroundings; think
before interacting with unusual objects

\textbf{Response:} Students should be aware of their vaccination history
and how to proceed if they require a tetanus shot due to injury

\subsection{7.5 Toxic Waste Exposure}\label{toxic-waste-exposure}

Exposure may occur via ingestion, inhalation, or dermal contact and can
be unpredictable.

\textbf{Prevention:} Wear appropriate PPE; avoid contact with unknown
substances

\textbf{Response:} Remove from exposure; decontaminate; seek medical
attention

\section{8. Biotic Hazards}\label{biotic-hazards}

The field is home to a vast diversity of organisms. Being respectful of
their environment is crucial to their safety and our own.

\subsection{8.1 Allergies}\label{allergies}

\begin{tcolorbox}[enhanced jigsaw, colback=white, titlerule=0mm, leftrule=.75mm, left=2mm, toptitle=1mm, rightrule=.15mm, title=\textcolor{quarto-callout-important-color}{\faExclamation}\hspace{0.5em}{Important}, opacitybacktitle=0.6, bottomrule=.15mm, breakable, opacityback=0, colframe=quarto-callout-important-color-frame, coltitle=black, arc=.35mm, bottomtitle=1mm, colbacktitle=quarto-callout-important-color!10!white, toprule=.15mm]

All allergies must be recorded in medical files before the start of
field work. Researchers with life-threatening allergies must bring their
Epi-pen to the field and keep it accessible at all times.

\end{tcolorbox}

\textbf{Mild allergic reaction:} Administer 25--50 mg of Diphenhydramine
(Benadryl)

\textbf{Severe allergic reaction} (swelling of face/throat, respiratory
distress): Administer 0.3 mg of Epinephrine; may re-administer after 5
minutes. Follow with 25--50 mg of Diphenhydramine and 20--40 mg of
Prednisone. Evacuate from field site immediately.

\subsection{8.2 Biotic Risk Summary}\label{biotic-risk-summary}

\begin{longtable}[]{@{}
  >{\raggedright\arraybackslash}p{(\linewidth - 6\tabcolsep) * \real{0.1923}}
  >{\raggedright\arraybackslash}p{(\linewidth - 6\tabcolsep) * \real{0.1923}}
  >{\raggedright\arraybackslash}p{(\linewidth - 6\tabcolsep) * \real{0.1346}}
  >{\raggedright\arraybackslash}p{(\linewidth - 6\tabcolsep) * \real{0.4808}}@{}}
\toprule\noalign{}
\begin{minipage}[b]{\linewidth}\raggedright
Organism
\end{minipage} & \begin{minipage}[b]{\linewidth}\raggedright
Habitats
\end{minipage} & \begin{minipage}[b]{\linewidth}\raggedright
Risks
\end{minipage} & \begin{minipage}[b]{\linewidth}\raggedright
Prevention and Response
\end{minipage} \\
\midrule\noalign{}
\endhead
\bottomrule\noalign{}
\endlastfoot
Mosquitos & Near stagnant water & Vectors for disease (e.g., Malaria) &
Use mosquito repellent (e.g., DEET products) \\
Rodents & Debris, dense underbrush, burrow holes & Disease, infection &
Do not touch rodents; if bitten, clean and disinfect \\
Mountain Lions & Native to North/South America; active early and late in
the day & Severe injury & Make yourself look larger; do not run; throw
rocks or sticks; avoid being alone \\
Snakes (Rattlesnakes, Cottonmouth) & Various & Snakebite & Walk in open
areas; wear heavy boots; listen for rattle; back away slowly; if bitten,
seek immediate medical attention \\
Spiders (Black Widow, Brown Recluse) & Dark sheltered areas & Spider
bite, nausea & Wear gloves while working; shake out clothing and
bedding; go to hospital if bitten \\
Bees and Wasps & Various & Swelling, pain, allergic reaction,
anaphylactic shock & Avoid disturbing nests; stay calm when pursued;
administer epinephrine for severe reactions; take to hospital \\
Fleas and Ticks & Underbrush, wooded areas & Lyme disease & Tick check
after fieldwork; use insect repellent; wear long clothing; suffocate
tick with petroleum jelly before removing with tweezers \\
Poison Oak & Riparian habitats & Itchy rash; red, swollen skin & Learn
to identify; apply Tecnu beforehand; wash affected areas with dish
soap \\
Stinging Nettle & Riparian habitats, meadows & Stinging sensation &
Learn to identify and avoid; use anti-itch cream if needed \\
\end{longtable}

\section{9. Other Illness and Injury}\label{other-illness-and-injury}

\subsection{9.1 Fractures and Breaks}\label{fractures-and-breaks}

Stabilize the injured area and avoid movement. Go to the nearest medical
facility for further care.

\subsection{9.2 Injury from Equipment}\label{injury-from-equipment}

Handle all equipment with caution. Carefully read operating procedures.
Seek help from fellow students or your instructor if you have questions
or doubts. Do not use any equipment if you are unsure of how to handle
it.

\subsection{9.3 Chronic Conditions (Epilepsy, Diabetes, Heart
Conditions)}\label{chronic-conditions-epilepsy-diabetes-heart-conditions}

Refer to emergency procedures. Report any chronic conditions to your
supervisor. Be aware of your surroundings and be prepared to react in an
emergency. Know where the first aid kit is and how to get to the nearest
medical facility.

\section{10. Field Preparation
Checklist}\label{field-preparation-checklist}

\subsection{10.1 Before Departure}\label{before-departure}

\begin{itemize}
\tightlist
\item[$\square$]
  Check weather report and air quality
\item[$\square$]
  Read and understand emergency procedures
\item[$\square$]
  Know locations of emergency equipment
\item[$\square$]
  Be aware of any special medical conditions of team members
\item[$\square$]
  Report to supervisor
\item[$\square$]
  Identify hospital closest to the field site
\item[$\square$]
  Develop a transportation plan with designated driver
\item[$\square$]
  Create a safety plan (itinerary, emergency contacts, risk assessment)
\end{itemize}

\subsection{10.2 Supply Checklist}\label{supply-checklist}

\begin{itemize}
\tightlist
\item[$\square$]
  Personal protective equipment
\item[$\square$]
  Rain gear (if necessary)
\item[$\square$]
  No excessively loose clothing; no sandals or open-toed shoes
\item[$\square$]
  Additional clothing/gear for environmental conditions (e.g., poison
  oak areas)
\item[$\square$]
  Sunscreen and sun hat (if necessary)
\item[$\square$]
  Full water bottle(s)
\item[$\square$]
  Rite in the Rain lab notebook and writing utensil
\item[$\square$]
  Watch
\item[$\square$]
  First aid kit
\item[$\square$]
  Cell phone with the following numbers programmed:

  \begin{itemize}
  \tightlist
  \item
    Emergency personnel
  \item
    Lab supervisor and/or instructor
  \item
    Lab teammates
  \end{itemize}
\end{itemize}

\subsection{10.3 On Site}\label{on-site}

\begin{itemize}
\tightlist
\item
  Identify a time to meet up and depart (if splitting up)
\item
  Leave no trace: ensure that the environment remains undisturbed
\item
  Supervisor ensures that safety procedures are followed
\item
  Report any potential field hazards to supervisor
\end{itemize}

\section{11. Emergency Response
Procedure}\label{emergency-response-procedure}

\subsection{11.1 Responding to a Medical
Emergency}\label{responding-to-a-medical-emergency}

\begin{enumerate}
\def\labelenumi{\arabic{enumi}.}
\tightlist
\item
  Survey the scene
\item
  Identify any potential risks that could affect other team members
\item
  Do not move the injured person unless necessary
\item
  Summon medical help and share the following information:

  \begin{itemize}
  \tightlist
  \item
    Suspected type of injury or illness
  \item
    Location
  \item
    Type of assistance required
  \end{itemize}
\item
  Identify a location for entry if an emergency vehicle is being
  summoned
\end{enumerate}

\subsection{11.2 Documenting the
Incident}\label{documenting-the-incident}

\begin{itemize}
\tightlist
\item
  Identify what happened and how the injury or illness occurred
\item
  This information will be used to eliminate hazards and prevent future
  incidents
\item
  Report the injury or illness to the lab supervisor
\end{itemize}

\section{12. References}\label{references}

\begin{itemize}
\tightlist
\item
  APHA, AWWA, WEF. (2012). \emph{Standard Methods for Examination of
  Water and Wastewater}. 22nd ed.~American Public Health Association
  (Eds.). Washington. 1360 pp.
\end{itemize}

\section{13. Field Work Planning Form}\label{field-work-planning-form}

Complete this form before each field outing:

\begin{itemize}
\tightlist
\item
  \textbf{Phone Numbers:} \_\_\_\_\_\_\_\_\_\_\_\_\_\_\_
\item
  \textbf{Location:} \_\_\_\_\_\_\_\_\_\_\_\_\_\_\_
\item
  \textbf{Risk Assessment:} \_\_\_\_\_\_\_\_\_\_\_\_\_\_\_
\item
  \textbf{Risk Mitigation Plan:} \_\_\_\_\_\_\_\_\_\_\_\_\_\_\_
\item
  \textbf{Emergency Services:} \_\_\_\_\_\_\_\_\_\_\_\_\_\_\_
\end{itemize}

\section{14. Training Record}\label{training-record}

Training is documented in EA Program files. Before working
independently, personnel must:

\begin{enumerate}
\def\labelenumi{\arabic{enumi}.}
\tightlist
\item
  Read this SOP in full
\item
  Attend a field safety orientation
\item
  Demonstrate understanding of emergency procedures
\item
  Be signed off by the instructor or designated trainer
\end{enumerate}




\end{document}
